% META: {"id": "TCOMPL-INEG-EX-007", "chapitre": "TCOMPL-INEGALITES", "type_objet": "exercice", "capacites": ["TCOMPL-INEG-C5"], "capacites_codes": ["C5"], "methodes": ["M5"], "parcours": 3, "competences": ["calculer", "raisonner"], "duree_min": 18, "mode_creation": "ex_nihilo", "sources_inspiration": [], "fichier_tex": "chapitres/TCOMPL-INEGALITES/exercices/TCOMPL-INEG-EX-007.tex", "corrige_tex": "chapitres/TCOMPL-INEGALITES/corriges/TCOMPL-INEG-CO-007.tex", "status": "generated"}
% BEGIN-VERIFY
% from sympy import symbols, integrate, Rational
% x = symbols('x')
% aire3 = integrate(x - x ** 3, (x, 0, 1))
% assert aire3 == Rational(1, 4)
% gini3 = 2 * aire3
% assert gini3 == Rational(1, 2)
% aire2 = integrate(x - x ** 2, (x, 0, 1))
% assert aire2 == Rational(1, 6)
% gini2 = 2 * aire2
% assert gini2 == Rational(1, 3)
% assert gini3 > gini2
% assert 0 <= gini2 <= 1 and 0 <= gini3 <= 1
% END-VERIFY
\begin{exercice}{TCOMPL-INEG-EX-007}{3}{18}
On garde $L(x)=x^3$.
\begin{enumerate}
  \item Calculer l'aire entre la première bissectrice et la courbe de Lorenz
        sur $[0\,;\,1]$.
  \item En déduire l'indice de Gini, égal au double de cette aire.
  \item Comparer avec l'indice obtenu pour $L(x)=x^2$. Quelle répartition est
        la plus inégalitaire ?
\end{enumerate}
\end{exercice}
